\section{Relation to earlier theories and scope}\label{sec:literature}
The static barycentric step belongs to the theory of comparison of experiments and convex order. Blackwell's comparison theorem\cite{blackwell} concerns the decision content of a parameter-independent garbling; Strassen's theorem\cite{strassen} identifies convex order with a martingale coupling. Lemma~\ref{lem:allocation} is a finite-target proof of this classical mechanism. Bayesian persuasion\cite{kamenica} optimizes information structures subject to Bayes plausibility; coarse information design\cite{lyu} relates conditional-mean cells to the curvature of their value function. Here the incoming law is fixed by the actually executed physical acquisition. One may only garble that law, and the next law depends on the retained information and chosen action. Neither Bayes plausibility alone nor a static posterior partition enforces these serial constraints.

Causal transport\cite{backhoff} imposes nonanticipation on couplings and has its own dynamic programming principle. The present restriction is different from nonanticipation alone: retained information must have prescribed finite cardinalities, and a reused internal label must invoke one and the same report kernel at every phase. Equation~\eqref{eq:simdual} imposes this latter constraint simultaneously, rather than selecting an arbitrary adapted coupling at each time. The theorem does not identify directed information, a channel rate, or a transport cost with persistent label cardinality. The compact finite-obstruction result is a topological separation statement; it is not an efficient optimization algorithm over causal transport plans.

Several approximation theories quantify finite-memory performance for known POMDPs. Kara and Y\"uksel\cite{kara} construct finite-window approximations under filter stability and finite observation/action alphabets. Demirci, Kara and Y\"uksel\cite{refined} refine such estimates using expected Wasserstein or pathwise total-variation stability. Saldi, Y\"uksel and Linder\cite{saldi} establish finite-model approximation for discounted partially observed control under continuity conditions. These approaches compare approximate policies or values with their control benchmarks. They do not impose the same internally counted exact-deadline interface used in \eqref{eq:regular_aut}. Our quantitative theorem instead combines an actual-law converse with explicitly allocated finite observer states, while the exact allocation theorem itself does not require their stability conditions. This distinction does not imply that the present theorem solves the discounted or average-cost problems treated in those works.

Grover and Dimitrakakis\cite{grover} adapt belief discretization to depth in POMDP planning. Li, Hammar and Bertsekas\cite{li} aggregate beliefs through features and conditional means and derive value-approximation bounds. Adaptive grid size, feature aggregation and conditional means are therefore not new ingredients here. The additional requirement is calibration under all incoming edges of the deployed observer, and for an internally timed observer the compatibility of the same transition programme across all cuts. The mass--curvature estimate and block allocation used in Section~\ref{sec:geometry} are inherited quantitative mechanisms\cite{prior}; they are not counted again as new results merely because they now follow an exact feasibility theorem.

The recent belief-metric covering analysis of Zhu and Lu\cite{zhu} addresses off-policy evaluation and coverage/sample-efficiency bounds for stable policies in offline POMDP learning. Its geometric coverage condition is not the hard cardinality of a causal observer's persistent state. Kim\cite{kim} studies finite-memory belief mismatch and policy-conditional performance along a common closed-loop execution. That information pattern differs both from optimizing over every legal controller and from charging a stationary machine's timing states. These comparisons are at the stated theorem objectives; neither paper is dismissed on the basis of its terminology or publication date. In particular, this article does not claim matching acquisition complexity for learning an unknown raw kernel from offline data.

Zero-delay source coding is another close but distinct theory. Wood, Linder and Y\"uksel\cite{wood} and Ghomi, Linder and Y\"uksel\cite{ghomi} study optimal stationary codes and finite-memory approximations for Markov sources under long-run distortion criteria. Stavrou, Charalambous, Charalambous and Loyka\cite{stavrou} develop nonanticipative rate-distortion bounds and causal filtering constructions, including Gaussian processes. An information-rate or communication-alphabet constraint does not bound every encoder/decoder internal state. Conversely, our observer need not transmit one symbol over a communication channel at every physical call. Thus their infinite-horizon and information-rate statements cannot be substituted for the finite-label cut converse here, and our results do not subsume theirs.

Positive quantum instruments have a standard operational basis\cite{davies}. Theorem~\ref{thm:quantum} uses a deliberately explicit finite-dimensional positive preparation and spanning instrument, so that quotient separation, report continuity and acquired density can be verified without commutativity. It is not a theorem about arbitrary quantum memory, operator domains, or physically unrestricted measurement design. Likewise, Proposition~\ref{prop:simulation} is the usual bounded-risk comparison mechanism\cite{lecam} with the missing causal and resource interfaces made explicit, rather than a new total-variation inequality. The effective finite compilation theorem uses classical real algebraic decision methods\cite{bpr}; it gives no new complexity bound for quantifier elimination.

The new proposed mathematical contribution is the exact finite-support causal allocation frontier, its phase-compatible dual, the compact global finite-obstruction consequence, and their combination with explicitly counted implementation and distinct matching regimes. The zero-loss criterion is independent of full task curvature; the block-stable law and blind-transport law exhibit when, and why, quantitative conclusions differ. Full predictive quotients of arbitrary Borel experiments, a uniformly effective theory of all stationary compact-carrier optimizers, unknown-kernel minimax learning, and a matched optimal region for programme size, scratch space and physical time are not established. Neither the fixed-horizon theory nor its bounded-stopping error estimate is an infinite-horizon theorem.
